% The structure under discussion: a binary tree whose internal nodes hold the
% branching point of the division that made them, and whose leaves are cells.
% A cell's address is the sequence of turns that reaches its leaf.
\begin{figure}[htbp]
  \centering
  \begin{tikzpicture}[scale=1.0, level distance=13mm,
    level 1/.style={sibling distance=44mm},
    level 2/.style={sibling distance=22mm},
    level 3/.style={sibling distance=12mm}]
    \node[tree node] (root) {branching point $(1,2)$}
      child {node[tree node] {$(0,1)$}
        child {node[tree leaf] {\textsf{LL}}}
        child {node[tree leaf] {\textsf{LR}}}
      }
      child {node[tree node] {$(2,3)$}
        child {node[tree leaf] {\textsf{RL}}}
        child {node[tree node] {$(3,1)$}
          child {node[tree leaf] {\textsf{RRL}}}
          child {node[tree leaf] {\textsf{RRR}}}
        }
      };
    \node[font=\scriptsize, black!60, right=3mm of root] {the root region};
  \end{tikzpicture}
  \caption{The proposed structure. Each internal node stores the branching point
  of the division that created it --- the cell its single door sits on --- and
  each leaf is a region of the maze. A cell's \keyterm{address} is the sequence
  of turns from the root to its leaf, so \textsf{RRL} names one region in three
  characters. Sibling leaves are the two halves a single division separated, and
  are joined in the maze by that division's door alone.}
  \label{fig:addresses}
\end{figure}
